\section{Definition and Threat Model}

A purpose contract is the immutable tuple
$C=(i,p,D,\Pi,W,M,K,R,B,L,E)$: identity/version, purpose, dataset binding,
projection, workload, model, capabilities, release predicate, privacy budget,
lifecycle, and evidence contract. Execution starts only if the current approved
contract and every supplied binding verify. Every external effect is mediated
by $K$, and output remains quarantined until release, privacy settlement, and
evidence completion succeed.

For approved Task A, confidential $Z$ should add utility beyond public $X$. For
prohibited Task B, paired worlds keep $X$ and ground truth fixed while changing
$Z$; influence appears when released actions or scores differ. We estimate this
property over registered samples rather than claiming universal noninterference.

The non-TEE protocol includes malicious callers, adversarial model/workload
behavior, substitutions, route fallback, unauthorized capabilities, release
violations, privacy-budget attacks, and evidence tampering. It trusts the host
kernel, administrator, container runtime, and cryptographic/verifier
implementations. Host compromise, physical and hardware side channels, covert
channels outside mediated I/O, and legal compliance are out of scope.
